% !TeX root = ../main.tex
%TODO: AUTHOR do this experimento for next-region
\chapter{A Parameter-Count Control for Next-Category Prediction}
\label{apx:parameter-count-control}

Chapter~\ref{ch:mobiwac} attributes the next-category results to the check-in-level
representation rather than to the size of the model that reads it. This experiment tests the
competing explanation. The dedicated category model was made wider until it had approximately the
same number of trainable parameters as the complete joint model, holding the representation, the
folds, and the training protocol fixed. If capacity were what the category task responds to, the
wider model should improve on the narrower one. It does not.

The control covers Alabama under the four-seed protocol and California at one seed, the second
limited by a per-fold cost thirty-six times higher there. It is an additional dissertation analysis
and does not appear in the Chapter~\ref{ch:mobiwac} manuscript.
Section~\ref{sec:conc:answer} uses it to test one alternative explanation for the
final result.

\section{Result at a glance}
\label{sec:apx-parameter-summary}

Table~\ref{tab:apx-parameter-comparison} places the three results side by side. At Alabama, where
the control ran the full four-seed protocol, the wider dedicated model does not improve on the
narrower one it replaces: it scores $0.53$ macro-F1 points lower with $6.5$ times the parameters,
and a paired test over the four seeds separates that difference from zero ($p = 0.0011$). The
direction is unanimous at the fold level, all twenty folds included, which is a stronger statement
than the seed-level mean alone. The joint
model, which carries a comparable budget to the wider one, sits between them and above the wider
model by $0.35$ points ($p = 0.041$). Capacity is therefore not the variable next-category
prediction responds to at this dataset, which leaves the representation as the explanation
Chapter~\ref{ch:mobiwac} gives.

At California the same three arms were measured at one seed, where the three lie within $0.06$
points of each other. That is a single-seed reading rather than a result under the four-seed
convention, and it is reported as such: no test is applied to it and no average is taken across the
two datasets.

\begin{table}[H]
    \centering
    \caption{Next-category macro-F1 for the original dedicated, wider dedicated,
    and joint models. The last two columns report wider minus original and joint
    minus wider, in macro-F1 points. Alabama: four seeds $\times$ five folds per arm;
    $\pm$ is the standard deviation across seeds. $^{\dagger}$California: seed 0 only,
    five folds, for all three arms; a single-seed reading, not a result under the
    four-seed convention.}
    \label{tab:apx-parameter-comparison}
    {\small
    \setlength{\tabcolsep}{4pt}
    \begin{tabular}{lrrrrr}
        \toprule
        Dataset & Original & Wider & Joint & Wider $-$ original & Joint $-$ wider \\
        \midrule
        Alabama                  & $30.77 \pm 0.07$ & $30.24 \pm 0.13$ & $30.59 \pm 0.07$ & $-0.53$ & $+0.35$ \\
        California$^{\dagger}$   & $35.62$           & $35.56$           & $35.61$           & $-0.06$ & $+0.05$ \\
        \bottomrule
    \end{tabular}
    }
\end{table}

The joint column uses the diagnostic-best category checkpoint, matching the
per-task convention of this control. Chapter~\ref{ch:mobiwac} instead uses one
jointly selected checkpoint for its headline table. Those values are not mixed here.

\section{Experimental design}
\label{sec:apx-parameter-design}

The control changes one factor: the hidden dimension of the dedicated
next-category model. At Alabama each arm uses four seeds, \{0, 1, 7, 100\}, and five
user-disjoint folds per seed, for 20 fitted models per arm, which is the convention the rest of
this document uses. At California one fold of the wider model takes $3{,}682$ seconds against
$102$ seconds at Alabama, so the four-seed protocol would take about twenty hours there; the
wider arm was measured at seed 0 and stopped. The reported value is macro-F1 under the same
epoch-selection rule as the arms it is compared against.

Table~\ref{tab:apx-parameter-counts} shows the capacity match. The original model
uses a hidden dimension of 256. The wider models reach 100.2\% and 101.9\% of the
joint-model parameter counts, so they do not receive a smaller parameter budget. Relative to the
original dedicated model, the wider models carry $6.5$ and $8.1$ times the trainable parameters.

\begin{table}[H]
    \centering
    \caption{Trainable parameter counts in the parameter-count control.}
    \label{tab:apx-parameter-counts}
    \begin{tabular}{lrrr}
        \toprule
        Dataset & Joint model & Original dedicated & Wider dedicated \\
        \midrule
        Alabama    & 4,197,621 & 644,359 & 4,207,399 ($h=672$) \\
        California & 5,151,189 & 644,359 & 5,249,719 ($h=752$) \\
        \bottomrule
    \end{tabular}
\end{table}

Because width can change the best learning rate, the wider model was given its own small search
before the reported runs, and it selects 0.0025 rather than the 0.005 the original model uses at
both datasets. Table~\ref{tab:apx-parameter-arms} reports the selected setting per dataset. This
search is smaller than the one the original dedicated model received, which is a reason the wider
arm is not read as a tuned ceiling.

\begin{table}[H]
    \centering
    \caption{The selected training setting for the wider dedicated category model at each dataset,
    with the macro-F1 it reaches. Alabama: four seeds $\times$ five folds, $\pm$ across seeds.
    $^{\dagger}$California: seed 0, five folds.}
    \label{tab:apx-parameter-arms}
    \begin{tabular}{lrrrc}
        \toprule
        Dataset & Hidden dimension & Batch size & Learning rate & Macro-F1 \\
        \midrule
        Alabama                 & 672 & 2,048 & 0.0025 & $30.24 \pm 0.13$ \\
        California$^{\dagger}$  & 752 & 8,192 & 0.0025 & $35.56$ \\
        \bottomrule
    \end{tabular}
\end{table}

\section{Interpretation and limits}
\label{sec:apx-parameter-result}

At Alabama, parameter count does not account for next-category performance: multiplying the
dedicated model's trainable parameters by $6.5$ lowers its macro-F1. The competing explanation this
control was built to test is therefore not supported, which leaves the input representation as the
factor Chapter~\ref{ch:mobiwac} identifies. What the control does not do is decompose the joint
model: it holds the representation fixed and varies width, so it says nothing about which of the
joint model's own mechanisms contribute to its results.

The scope is narrow and worth stating plainly: next-category prediction only, one wider model per
dataset, width rather than depth, a smaller hyperparameter search than the original model received,
the full four-seed protocol at Alabama and a single seed at California. The result rules out a
simple parameter-count explanation at Alabama and is consistent with the same reading at
California, but it is neither a complete capacity study nor a component ablation.

\noindent\textbf{Result files.} The per-fold scores of every arm, the driver log with per-seed
wall-clock times, and the parsed digests are tracked under
\path{docs/results/closing_data/apxi_v18/}.
